\documentclass[10pt,a4paper]{amsart}
\usepackage[margin=19mm]{geometry}
\usepackage{amsmath,amssymb,mathtools}
\usepackage[scr=rsfs,cal=euler]{mathalfa}
\usepackage{xfrac}
\usepackage[unicode,hidelinks,bookmarks=false]{hyperref}
\newcommand{\ie}{i.e.\ }
\newcommand{\cf}{cf.\ }
\newcommand{\resp}{respectively\ }
\newcommand{\Subsection}{\subsection}
\providecommand{\nobreakdash}{\nobreak\hbox{-}}
\setlength{\emergencystretch}{2em}
\multlinegap0pt
\usepackage{bm}
\usepackage{bbm}
\usepackage{dsfont}
\usepackage{xspace}
\usepackage{cancel}
\usepackage{mathtools}
\usepackage{cleveref}
\usepackage{amsthm}
\makeatletter
\@ifundefined{theorem}{\newtheorem{theorem}{Theorem}}{}
\makeatother
\title[Rational Neural Networks have Expressivity Advantages]{Rational Neural Networks have Expressivity Advantages}
\author{Maosen Tang, Alex Townsend}
\date{Source: arXiv:2602.12390v2 (CC BY 4.0)}
\begin{document}
\maketitle

\noindent\emph{Source and licence.} Rational Neural Networks have Expressivity Advantages. Version arXiv:2602.12390v2 (CC BY 4.0), distributed under the Creative Commons Attribution 4.0 International licence, as recorded in the arXiv licence metadata for this exact version and shown on the abstract page https://arxiv.org/abs/2602.12390v2. Full rights notice: licenses/arxiv-2602.12390-CC-BY-4.0.txt.\par\smallskip

\noindent\textit{Editorial scope.} Counted unit: the Theorem whose source label is \texttt{None} in the article (source0.tex lines 1511--1513, source label \texttt{None}), delivered together with the article's own complete original proof of it (source0.tex lines 1515--1656, one \texttt{proof} environment of 142 lines). No mathematics is written, repaired or re-derived: statement and proof are the article's own text line for line. Required context retained uncounted: none. Every definition, display and label of those lines is retained unchanged. Omitted from this package and not redistributed: 1--1510, 1514--1514, 1657--2169. Nothing inside a retained excerpt is omitted or altered: every retained body is delivered exactly as the article's lines are written, so no omission receipt of any kind accompanies this package. Citations: the retained excerpts carry no citation command at all, so no bibliography entry of the article is transcribed and no reference is redistributed. Reference adaptation: none was needed and none was made; every reference inside the delivered unit resolves inside it, so no unresolved reference or citation is produced. Printed slips kept literal: no printed word, symbol, hypothesis or number of the counted statement or of its proof was altered. Layout and class: the article's own class is \texttt{article} with packages including placeins, microtype, adjustbox, graphicx, multirow, subfigure, booktabs, comment, tabularx, siunitx, xcolor, hyperref, icml2026, amsmath; the wrapper is the lane's pinned \texttt{amsart} editorial wrapper at 10pt on A4, which restates the theorem-like environments the retained text needs and adds the article's own macro definitions that the retained text needs (none), with extra wrapper packages (bm, bbm, dsfont, xspace, cancel, mathtools, cleveref). The delivered numbering is the unit's own. Assets: no retained span contains a figure environment, a graphics-inclusion command, a PDF-page inclusion, a table or a TikZ picture; no image file is copied into this package, the package declares no asset dependency and no image file is redistributed with it. Licence: version arXiv:2602.12390v2 of this article is distributed under the Creative Commons Attribution 4.0 International licence, as recorded in the arXiv licence metadata for this exact version and shown on the abstract page https://arxiv.org/abs/2602.12390v2. A full rights notice is delivered at licenses/arxiv-2602.12390-CC-BY-4.0.txt. Characters: every retained excerpt is pure ASCII, so the delivered engine, XeLaTeX with its default Latin Modern text font, reports no missing character.
\begin{theorem}
 {Let $0<\varepsilon<\tfrac{1}{8}$. There exists a rational function $R_\varepsilon:[-1,1]\to[0,1]$ such that for any scalar-input, scalar-output GELU network, $F$, with uniformly bounded weights and biases that satisfies $\|F-R_\varepsilon\|_{L^\infty([-1,1])}\le \varepsilon$, the size of the network is $\Omega\bigl(\log(1/\varepsilon)\bigr)$.}
\end{theorem}

\begin{proof}
 {Let $\eta=\sqrt{\varepsilon}$ and consider
\[
R_\varepsilon(x)=\frac{\eta^2}{x^2+\eta^2}.
\]
Then $0<R_\varepsilon(x)\le 1$ on $[-1,1]$, $R_\varepsilon\in C^\infty([-1,1])$, and $R_\varepsilon''(0)=-2/\eta^2=-2/\varepsilon$. We quantify the curvature required of any smooth approximation to $R_\varepsilon$ at accuracy $\varepsilon$. Suppose $f\in C^2([-1,1])$ satisfies $\|f-R_\varepsilon\|_{L^\infty([-1,1])}\le \varepsilon$. Then for any fixed $h\in(0,1]$, a standard finite-difference form of Taylor's theorem yields a $\xi\in(-h,h)$ such that}
\begin{equation}
f(h)+f(-h)-2f(0) \;=\; f''(\xi)\,h^2,\qquad |f''(\xi)| \;=\; \frac{|f(h)+f(-h)-2f(0)|}{h^2}.
\label{eq:fpp-fd}
\end{equation}
 {Since $R_\varepsilon$ is even, $R_\varepsilon(h)=R_\varepsilon(-h)$, and}
\[
R_\varepsilon(h)+R_\varepsilon(-h)-2R_\varepsilon(0)
=
\frac{2\eta^2}{h^2+\eta^2}-2
=
-\frac{2h^2}{h^2+\eta^2}.
\]
 {Using $\|f-R_\varepsilon\|_\infty\le \varepsilon$ and the triangle inequality,}
\[
|f(h)+f(-h)-2f(0)|
\ge |R_\varepsilon(h)+R_\varepsilon(-h)-2R_\varepsilon(0)|-4\varepsilon
=
\frac{2h^2}{h^2+\eta^2}-4\varepsilon.
\]
Combining with~\cref{eq:fpp-fd} gives
\begin{equation}
|f''(\xi)|
\ge \frac{2}{h^2+\eta^2}-\frac{4\varepsilon}{h^2}.
\label{eq:curv-lb-h}
\end{equation}
Select $h=\eta$ (which lies in $(0,1]$ since $\eta=\sqrt{\varepsilon}<1$). Then~\cref{eq:curv-lb-h} becomes
\[
|f''(\xi)|
\ge \frac{1}{\eta^2}-\frac{4\varepsilon}{\eta^2}
=
\frac{1-4\varepsilon}{\eta^2}.
\]
Because $0<\varepsilon<\tfrac18$, we have $1-4\varepsilon>\tfrac12$, hence
\[
|f''(\xi)|\ge \frac{1}{2\eta^2}=\frac{1}{2\varepsilon}.
\]
 {In particular, if $F$ is any GELU network with $\|F-R_\varepsilon\|_\infty\le \varepsilon$, then $F\in C^\infty$ and the above applies with $f=F$, giving}
\begin{equation}
\sup_{x\in[-1,1]} |F''(x)| \;\ge\; \frac{1}{2\varepsilon}.
\label{eq:need-curv}
\end{equation}

Next, we upper bound the curvature available to a bounded-parameter GELU network. Consider a depth-$L$ fully connected scalar-input, scalar-output network with hidden widths $W_1,\dots,W_L\in\mathbb N$:
\[
h^{(0)}(x)=x,\quad z^{(\ell)}=A^{(\ell)}h^{(\ell-1)}+b^{(\ell)},\quad h^{(\ell)}=G(z^{(\ell)}),\quad F(x)=c^\top h^{(L)}+d,
\]
where $A^{(\ell)}\in\mathbb R^{W_\ell\times W_{\ell-1}}$ with $W_0=1$, $b^{(\ell)}\in\mathbb R^{W_\ell}$, $c\in\mathbb R^{W_L}$, and all entries of $A^{(\ell)},b^{(\ell)},c,d$ have magnitude at most $B\ge1$, where $B$ is independent of $\varepsilon$. Moreover, $G$ is applied entrywise to vector inputs. For
\[
G(z)=\tfrac{z}{2}\bigl(1+\tanh(\alpha(z+\beta z^3))\bigr),\qquad \alpha=\sqrt{2/\pi},\ \beta=0.044715,
\]
one has finite constants
\[
C_1:=\sup_{z\in\mathbb R}|G'(z)|<\infty,\qquad C_2:=\sup_{z\in\mathbb R}|G''(z)|<\infty.
\]
Set \(
C:=\max\{1,\,C_1,\,\sqrt{C_2}\},
\) so that $|G'(z)|\le C$ and $|G''(z)|\le C^2$ for all $z$.

With the mixed operator norm, for any matrix $M$ write $\|M\|_{1\to\infty}=\max_i\sum_j|M_{ij}|$; in particular $\|A^{(\ell)}\|_{1\to\infty}\le B\,W_{\ell-1}$ and $\|c\|_1\le B\,W_L$ (where $\|c\|_1=\sum_i |c_i|$). To propagate derivatives through the layers we track the suprema
\[
S_\ell=\sup_{x\in[-1,1],\,1\le i\le W_\ell}\Bigl|\frac{\partial h^{(\ell)}_i}{\partial x}\Bigr|,\qquad
T_\ell=\sup_{x\in[-1,1],\,1\le i\le W_\ell}\Bigl|\frac{\partial^2 h^{(\ell)}_i}{\partial x^2}\Bigr|.
\]
By the chain rule, for each $\ell\ge1$ and each coordinate $i$,
\[
\Bigl|\frac{d}{dx}h^{(\ell)}_i(x)\Bigr|
=|G'(z^{(\ell)}_i(x))|\,\Bigl|\sum_j A^{(\ell)}_{ij}\,\frac{d}{dx}h^{(\ell-1)}_j(x)\Bigr|
\le C\,\|A^{(\ell)}\|_{1\to\infty}\,S_{\ell-1},
\]
hence
\[
S_\ell \le C\,\|A^{(\ell)}\|_{1\to\infty}\,S_{\ell-1}\ \le\ (C B W_{\ell-1})\,S_{\ell-1},\qquad S_0=1,\ T_0=0,
\]
so \(
S_\ell \le \prod_{k=1}^{\ell} (C B W_{k-1}).
\) For the second derivatives, using $h^{(\ell)}_i{}''=G''(z^{(\ell)}_i)\,(z^{(\ell)}_i{}')^2+G'(z^{(\ell)}_i)\,z^{(\ell)}_i{}''$ and
$|z^{(\ell)}_i{}'|\le \|A^{(\ell)}\|_{1\to\infty}S_{\ell-1}$, $|z^{(\ell)}_i{}''|\le \|A^{(\ell)}\|_{1\to\infty}T_{\ell-1}$, we obtain
\[
T_\ell \le C^2\,\|A^{(\ell)}\|_{1\to\infty}^2\,S_{\ell-1}^2 + C\,\|A^{(\ell)}\|_{1\to\infty}\,T_{\ell-1}
\ \le\ (C B W_{\ell-1})^2 S_{\ell-1}^2 + (C B W_{\ell-1}) T_{\ell-1}.
\]
Since $C B W_{\ell-1}\ge1$, an induction gives \(
T_\ell \le \ell\,\Bigl(\prod_{k=1}^{\ell} C B W_{k-1}\Bigr)^{\!2}.
\) The output linear map contributes $\|c\|_1$, hence
\[
\sup_{x\in[-1,1]}|F''(x)| \le \|c\|_1\,T_L \le B W_L \, L\,\Bigl(\prod_{k=1}^{L} C B W_{k-1}\Bigr)^{\!2}
= L\,(C B)^{2L}\,B\,W_L\,\prod_{k=1}^{L-1} W_k^{2}.
\]
Using $B\le C B$ and $W_L\le W_L^2$ (since $C,B,W_L\ge1$), this relaxes to
\begin{equation}\label{eq:Fpp}
\sup_{x\in[-1,1]}|F''(x)| \le L\,(C B)^{2L+1}\,\prod_{k=1}^{L} W_k^{2}.
\end{equation}

Comparing~\cref{eq:need-curv} and~\cref{eq:Fpp} gives
\begin{equation}\label{eq:key}
\frac{1}{2\varepsilon}\ \le\ \sup_{x\in[-1,1]}|F''(x)|\ \le\ L\,(C B)^{2L+1}\,\prod_{k=1}^{L} W_k^{2}.
\end{equation}
Let $\mathsf{A}:=\tfrac{1}{2\varepsilon}$. From~\cref{eq:key} we obtain
\[
\prod_{k=1}^{L} W_k^{2}\ \ge\ \frac{\mathsf{A}}{L\,(C B)^{2L+1}}
\qquad\Longrightarrow\qquad
\prod_{k=1}^{L} W_k\ \ge\ \Bigl(\frac{\mathsf{A}}{L\,(C B)^{2L+1}}\Bigr)^{\!1/2}.
\]
To convert a product constraint into parameters we pass to the count of hidden-to-hidden weights \(
S:=\sum_{\ell=1}^{L-1} W_\ell W_{\ell+1},
\) a subset of the total parameter count, so $\mathrm{size}\ge S$. Since $W_\ell\ge1$,
\[
\prod_{\ell=1}^{L-1} (W_\ell W_{\ell+1})
= W_1 W_L \prod_{k=2}^{L-1} W_k^{2}
\ \ge\ \prod_{k=1}^{L} W_k.
\]
By AM--GM on $\{W_\ell W_{\ell+1}\}_{\ell=1}^{L-1}$,
\[
\Bigl(\frac{S}{L-1}\Bigr)^{\!L-1}\ \ge\ \prod_{\ell=1}^{L-1} (W_\ell W_{\ell+1})
\ \ge\ \Bigl(\frac{\mathsf{A}}{L\,(C B)^{2L+1}}\Bigr)^{\!1/2}.
\]
Hence, for $L\ge2$,
\[
S\ \ge\ (L-1)\,\exp\!\Bigl(\frac{\log \mathsf{A}-(2L+1)\log(CB)-\log L}{2(L-1)}\Bigr).
\]
Writing $k=\log(CB)\ge0$ and $a=\log \mathsf{A}>0$ (since $\varepsilon<\tfrac18$ implies $\mathsf{A}>4$), and noting $\tfrac{2L+1}{2(L-1)}\le \tfrac52$ and $\tfrac{\log L}{2(L-1)}\le \tfrac12$ for $L\ge2$, we get
\[
S\ \ge\ e^{-(\tfrac52 k+\tfrac12)}\,(L-1)\,e^{a/(2(L-1))}.
\]
For any $t>0$, the function $t\mapsto t\,e^{a/(2t)}$ is minimized at $t=a/2$, yielding $t\,e^{a/(2t)}\ge \tfrac{a}{2}\,e$, so with $t=L-1$, \(
S\ \ge\ c_0\,\log \mathsf{A}
\) for a constant $c_0>0$ depending only on $B$ and $G$.

The single-hidden-layer case follows directly from~\cref{eq:key}: when $L=1$,~\cref{eq:Fpp} gives $\mathsf{A}\le (C B)^3 W_1^2$, and the total size is at least $2W_1$, so $\mathrm{size}\ge c_1\,\mathsf{A}^{1/2}\ge c_2\,\log \mathsf{A}$ (since $\mathsf{A}>4$ on $0<\varepsilon<\tfrac18$), for constants $c_1,c_2>0$ depending only on $B$ and $G$.

Therefore, in all cases,
\[
\mathrm{size}\ \ge\ c\,\log \mathsf{A}\ =\ \Omega\bigl(\log(1/\varepsilon)\bigr),
\]
with a constant $c>0$ depending only on $B$ and $G$.
\end{proof}

\end{document}
